\chapter{Regresión logística}

\section{Apuntes del Tema 4: Regresión Logística}

\subsection{1. Concepto y Función Sigmoide}
\begin{itemize}
	\item \textbf{Objetivo:} Clasificación binaria prediciendo $P(y=1|x) \in (0,1)$.
	\item \textbf{Sigmoide:} $\sigma(z) = \frac{1}{1 + e^{-z}}$, donde $z = x^\top w + b$.
\end{itemize}

\subsection{2. Función de Pérdida (BCE)}
\begin{itemize}
	\item \textbf{Binary Cross-Entropy:} $\mathcal{L} = - [ y \log(\hat{y}) + (1 - y) \log(1 - \hat{y}) ]$.
	\item \textbf{Propiedad:} Estrictamente convexa, garantiza convergencia al mínimo global.
	\item \textbf{Derivación Probabilística:} Equivalente a la Máxima Verosimilitud (MLE) bajo una distribución de Bernoulli.
\end{itemize}

\subsection{3. Regularización ($L_1$ y $L_2$)}
\begin{itemize}
	\item \textbf{Ridge ($L_2$):} Penaliza $w^2$ (esfera), reduce magnitud de pesos.
	\item \textbf{Lasso ($L_1$):} Penaliza $|w|$ (diamante), induce dispersión (\emph{sparsity}), lleva pesos a cero.
\end{itemize}

\subsection{4. Práctica}
\begin{lstlisting}[language=Python]
	# Pipeline de Scikit-Learn con escalado, poly y logistic
	pipeline = Pipeline([
	('scaler', StandardScaler()),
	('poly', PolynomialFeatures(degree=2)),
	('logistic', LogisticRegression(penalty='l2'))
	])
\end{lstlisting}

\section{Autoevaluación}
\begin{itemize}
	\item \textbf{P:} ¿Por qué la salida logística es una probabilidad? \textbf{R:} La sigmoide mapea el espacio real al intervalo $(0,1)$.
	\item \textbf{P:} ¿Riesgo de polinomio grado alto? \textbf{R:} Sobreajuste; requiere regularización.
\end{itemize}